\documentclass[11pt,a4paper]{article}
\makeatletter\def\input@path{{../../templates/}}\makeatother
\usepackage{avaliacao-poo}
\configurarAvaliacao{Checkpoint 1}{Programação Orientada a Objetos}
\validarCheckpointDuasPaginas
\begin{document}
\cabecalhoAvaliacao
\blocoQuadroRespostas
\cabecalhoCenario{Elevador}
Os arquivos \codigoInline{Elevador.java} e \codigoInline{Main.java} estão no mesmo pacote.
Considere exatamente o código abaixo.

\noindent\begin{minipage}[t]{.37\textwidth}
\vspace{0pt}
\begin{painelreferencia}{Elevador.java}
\begin{lstlisting}[style=cenariocodigo]
public class Elevador {
    int andarAtual = 0;
    boolean portaAberta = false;

    void subir() {
        andarAtual = andarAtual + 1;
    }
    void abrirPorta() {
        portaAberta = true;
    }
    void fecharPorta() {
        portaAberta = false;
    }
    int getAndarAtual() {
        return andarAtual;
    }
    boolean getPortaAberta() {
        return portaAberta;
    }
}
\end{lstlisting}
\end{painelreferencia}
\end{minipage}\hfill
\begin{minipage}[t]{.61\textwidth}
\vspace{0pt}
\begin{painelreferencia}{Main.java}
\begin{lstlisting}[style=cenariocodigo]
public class Main {
    public static void main(String[] args) {
        Elevador principal = new Elevador();
        Elevador painel = principal;
        Elevador servico = new Elevador();

        principal.subir();
        painel.abrirPorta();

        System.out.println(principal.getAndarAtual());
        System.out.println(principal.getPortaAberta());

        painel.fecharPorta();
        painel.subir();

        System.out.println(painel.getAndarAtual());
        System.out.println(painel.getPortaAberta());

        servico.subir();
        servico.abrirPorta();

        System.out.println(servico.getAndarAtual());
    }
}
\end{lstlisting}
\end{painelreferencia}
\end{minipage}
\vspace{-\espacoS}
\inicioQuestoes

\noindent\begin{minipage}[t]{.32\textwidth}
\questaoDireta{1}{Faça o teste de mesa. Para cada linha indicada, escreva exatamente o valor impresso.}{25 pontos}
\itensDuasColunas{
  \itemSimples{a}{Linha 10}
  \itemSimples{b}{Linha 11}
  \itemSimples{c}{Linha 16}
}{
  \itemSimples{d}{Linha 17}
  \itemSimples{e}{Linha 22}
}

\end{minipage}\hfill
\begin{minipage}[t]{.66\textwidth}
\questaoDireta{2}{A classe agora protege os campos e controla a subida. Considere somente o código abaixo e marque cada afirmação como V ou F.}{25 pontos}
Os getters, \codigoInline{abrirPorta()} e \codigoInline{fecharPorta()} continuam iguais aos do cenário.
\begin{painelEvolucao}{Trechos alterados em Elevador.java}
\begin{lstlisting}[style=evolucaocodigo]
private int andarAtual;
private boolean portaAberta;
int subir() {
    if (portaAberta == false) {
        if (andarAtual < 3) {
            andarAtual = andarAtual + 1;
        }
    }
    return andarAtual;
}
\end{lstlisting}
\end{painelEvolucao}

\itemVF{a}{Em \codigoInline{Main}, a instrução \codigoInline{principal.andarAtual = 2;} compila.}
\itemVF{b}{Porta fechada, andar 2: \codigoInline{int resultado = principal.subir();} deixa \codigoInline{resultado} e o elevador no andar 3.}
\itemVF{c}{Porta aberta, andar 1: \codigoInline{subir()} retorna 1 e preserva o andar 1.}
\itemVF{d}{Porta fechada, andar 3: duas chamadas de \codigoInline{subir()} retornam 3 e preservam o andar 3.}
\itemVF{e}{\codigoInline{Main} pode consultar o campo privado pelo método \codigoInline{getAndarAtual()}.}
\end{minipage}

\noindent\begin{minipage}[t]{.485\textwidth}
\questaoDireta{3}{O sistema também registra chamados de manutenção, com descrição, prioridade, situação e a ação de encerrar. Qual alternativa mantém juntos esses dados e essa ação?}{5 pontos}
\begin{alternativas}
\item Manter descrição, prioridade e situação em variáveis separadas de \codigoInline{Main} e encerrar o chamado ali.
\item Criar uma classe apenas para a descrição e manter prioridade, situação e encerramento em \codigoInline{Main}.
\item Criar \codigoInline{ChamadoManutencao} com descrição, prioridade, situação e um método para encerrar o chamado.
\item Criar uma classe para cada dado do chamado e realizar o encerramento em \codigoInline{Main}.
\item Manter os três dados em \codigoInline{Main} e passá-los a um método externo sempre que o chamado mudar.
\end{alternativas}
\questaoDireta{4}{Na instrução \codigoInline{Elevador principal = new Elevador();}, qual alternativa identifica corretamente a classe, a variável e o objeto criado?}{5 pontos}
\begin{alternativas}
\item \codigoInline{Elevador} é a variável; \codigoInline{principal} é a classe; \codigoInline{new} copia um objeto.
\item \codigoInline{Elevador} é a classe; \codigoInline{principal} guarda uma referência; \codigoInline{new Elevador()} cria o objeto.
\item \codigoInline{principal} é o objeto; \codigoInline{Elevador} guarda a referência usada por \codigoInline{new}.
\item \codigoInline{new Elevador()} é a variável; \codigoInline{principal} declara outra classe.
\item Os três nomes representam o mesmo papel e identificam diretamente o objeto.
\end{alternativas}
\questaoDireta{5}{Quando \codigoInline{abrirPorta()} é executado, o valor de \codigoInline{portaAberta} muda. Qual alternativa classifica corretamente o campo e o método?}{5 pontos}
\begin{alternativas}
\item \codigoInline{portaAberta} é estado; \codigoInline{abrirPorta()} é comportamento.
\item \codigoInline{portaAberta} é comportamento; \codigoInline{abrirPorta()} é estado.
\item Ambos são estados porque pertencem ao mesmo objeto.
\item Ambos são comportamentos porque podem ser acessados com ponto.
\item \codigoInline{portaAberta} é variável local; \codigoInline{abrirPorta()} é método de \codigoInline{Main}.
\end{alternativas}
\questaoDireta{6}{Queremos um método que receba um andar de destino inteiro e não retorne valor. Qual assinatura atende ao requisito?}{5 pontos}
\begin{alternativas}
\item \codigoInline{int irPara()}
\item \codigoInline{void irPara()}
\item \codigoInline{int irPara(int destino)}
\item \codigoInline{void irPara(int destino)}
\item \codigoInline{void irPara(float portaAberta)}
\end{alternativas}
\questaoDireta{7}{Observe o código. Ele compila? Se não, qual linha impede a compilação?}{5 pontos}
\begin{painelEvolucao}{Trecho em Elevador.java}
\begin{lstlisting}[style=questaocodigo]
int andarAtual;
void mostrar() {
    int anterior;
    System.out.println(andarAtual);
    System.out.println(anterior);
}
\end{lstlisting}
\end{painelEvolucao}
\begin{alternativas}
\item Compila e as linhas 4 e 5 imprimem \codigoInline{0}.
\item Não compila: a linha 5 tenta ler a variável local \codigoInline{anterior} sem inicializá-la.
\item Não compila: a linha 4 tenta ler o campo \codigoInline{andarAtual} sem inicializá-lo.
\item Compila, mas a execução falha na linha 5 ao acessar \codigoInline{anterior}.
\item Não compila: apenas declarar a variável local na linha 3 já é inválido.
\end{alternativas}
\end{minipage}\hfill
\begin{minipage}[t]{.485\textwidth}
\questaoDireta{8}{\codigoInline{Elevador} terá o campo \codigoInline{int quantidadePessoas}; a lotação máxima é 8. Onde deve ficar a decisão de aceitar ou rejeitar um grupo?}{5 pontos}
\begin{alternativas}
\item \codigoInline{Main} consulta a lotação e decide se inclui o grupo.
\item Um auxiliar de \codigoInline{Main} calcula e altera a lotação.
\item \codigoInline{Main} verifica o grupo e outro método externo verifica o limite.
\item Um setter recebe a nova lotação calculada fora do objeto.
\item \codigoInline{Elevador.embarcar(quantidade)} verifica sua lotação e o limite.
\end{alternativas}
\questaoDireta{9}{Depois de executar a última linha abaixo, quais são os valores de \codigoInline{diferenca} e \codigoInline{andarAtual}?}{5 pontos}
\begin{painelEvolucao}{Trecho em Elevador.java e Main.java}
\begin{lstlisting}[style=questaocodigo]
// Em Elevador.java
int calcularDiferenca(int destino) {
    return destino - andarAtual;
}
// Em Main.java
Elevador teste = new Elevador();
teste.subir();
int diferenca = teste.calcularDiferenca(3);
\end{lstlisting}
\end{painelEvolucao}
\begin{alternativas}
\item \codigoInline{diferenca} recebe \codigoInline{3} e \codigoInline{andarAtual} passa a ser \codigoInline{3}.
\item \codigoInline{diferenca} recebe \codigoInline{2} e \codigoInline{andarAtual} passa a ser \codigoInline{3}.
\item Não compila: um método com retorno \codigoInline{int} não pode receber parâmetro.
\item \codigoInline{diferenca} recebe \codigoInline{2} e \codigoInline{andarAtual} continua em \codigoInline{1}.
\item O método não retorna valor porque não foi declarado com \codigoInline{void}.
\end{alternativas}
\questaoDireta{10}{Após executar o trecho, o que é correto sobre os estados de \codigoInline{a} e \codigoInline{b} e sobre \codigoInline{a == b}?}{5 pontos}
\begin{painelEvolucao}{Trecho em Main.java}
\begin{lstlisting}[style=questaocodigo]
Elevador a = new Elevador();
Elevador b = new Elevador();
a.subir();
\end{lstlisting}
\end{painelEvolucao}
\begin{alternativas}
\item Antes da chamada, \codigoInline{a == b} é \codigoInline{true}; depois, ambos ficam no andar 1.
\item Antes da chamada, os estados diferem; depois, somente \codigoInline{b} fica no andar 1.
\item Antes da chamada, os estados são iguais e \codigoInline{a == b} é \codigoInline{false}; depois, somente \codigoInline{a} fica no andar 1.
\item Antes da chamada, \codigoInline{a == b} é \codigoInline{true}; depois, somente \codigoInline{a} fica no andar 1.
\item Antes da chamada, os estados são iguais; por isso, alterar \codigoInline{a} também altera \codigoInline{b}.
\end{alternativas}
\questaoDireta{11}{Uma classe de outro pacote tenta chamar \codigoInline{getAndarAtual()}, que não tem modificador de acesso. O que acontece?}{5 pontos}
\begin{alternativas}
\item Compila porque a classe \codigoInline{Elevador} é pública.
\item Não compila: o acesso de pacote permite uso apenas no mesmo pacote.
\item Só compila se o campo \codigoInline{andarAtual} também for público.
\item O método se torna privado automaticamente por ser chamado de outro pacote.
\item Compila porque a ausência de modificador equivale a \codigoInline{public}.
\end{alternativas}
\questaoDireta{12}{Os andares válidos vão de 0 a 3. Após acrescentar o método abaixo, \codigoInline{Main} chama \codigoInline{principal.definirAndar(8)}. O que acontece?}{5 pontos}
\begin{painelEvolucao}{Trecho acrescentado em Elevador.java}
\begin{lstlisting}[style=questaocodigo]
public void definirAndar(int novoAndar) {
    andarAtual = novoAndar;
}
\end{lstlisting}
\end{painelEvolucao}
\begin{alternativas}
\item Não compila: a própria classe não pode alterar um campo privado.
\item Não compila: o valor 8 está fora do intervalo permitido.
\item Compila, e \codigoInline{andarAtual} passa automaticamente para 3.
\item Compila, e \codigoInline{andarAtual} passa para 8 porque não há validação.
\item Compila, mas \codigoInline{andarAtual} continua em 0 porque o método é \codigoInline{void}.
\end{alternativas}
\end{minipage}

\end{document}
